\documentclass[10pt,reqno]{amsart}

\usepackage[letterpaper,margin=1in]{geometry}
\usepackage{amsmath,amssymb,mathtools,microtype}
\usepackage[hidelinks]{hyperref}
\hypersetup{pdftitle={Antichain Subset Sums in Rank-Three Multiplicative Semigroups}}

\setlength{\parindent}{1.15em}
\setlength{\parskip}{0pt}
\setlength{\abovedisplayskip}{5pt plus 2pt minus 2pt}
\setlength{\belowdisplayskip}{5pt plus 2pt minus 2pt}
\setlength{\abovedisplayshortskip}{3pt plus 1pt}
\setlength{\belowdisplayshortskip}{3pt plus 1pt}
\setlength{\jot}{1.5pt}

\newtheorem{theorem}{Theorem}[section]
\newtheorem{proposition}[theorem]{Proposition}
\newtheorem{lemma}[theorem]{Lemma}
\newtheorem{corollary}[theorem]{Corollary}
\newtheorem{remark}[theorem]{Remark}

\newcommand{\N}{\mathbb N}
\newcommand{\Nzero}{\mathbb Z_{\ge 0}}
\newcommand{\Z}{\mathbb Z}
\newcommand{\R}{\mathbb R}
\newcommand{\T}{\mathbb T}
\newcommand{\Sg}{\mathcal S}
\newcommand{\B}{\mathcal B}
\newcommand{\Eset}{\mathcal E}
\newcommand{\Prob}{\mathbb P}
\newcommand{\Ex}{\mathbb E}
\newcommand{\e}{\mathrm e}
\newcommand{\distz}[1]{\lVert #1\rVert}

\title[Antichain subset sums]{Antichain Subset Sums in Rank-Three Multiplicative Semigroups}
\author{}
\date{}

\begin{document}

\begin{abstract}
Let $a,b,c\ge2$ be pairwise coprime and
$\Sg=\{a^k b^\ell c^m:k,\ell,m\ge0\}$.  For every real
$1<\rho<\min(a,b,c)$, every sufficiently large integer is a subset sum of
$\Sg\cap[x,\rho x)$ for some $x$.  Such a set is automatically a divisibility
antichain, resolving Erd\H{o}s Problem~123.  For rational $\rho$ we prove the
stronger statement that the Bernoulli subset sums of each large band satisfy a
uniform local central limit theorem.  The main new ingredient is an arithmetic
rigidity argument on a triangular grid in exponent space; the remaining
Fourier analysis is standard.
\end{abstract}

\maketitle

\section{Statement and reduction}

Fix pairwise-coprime integers $a,b,c\ge2$, and write
\[
 \Sg=\{a^k b^\ell c^m:k,\ell,m\in\Nzero\},
 \qquad d=\min(a,b,c).
\]
Let $\rho=p/q\in\mathbb Q$ satisfy $1<\rho<d$, with $p,q\in\N$ and $q>0$.
For an integer $x\ge1$ set
\[
 \B_x=\Sg\cap[x,\rho x),\qquad
 W_x=\sum_{s\in\B_x}s,\qquad V_x^2=\sum_{s\in\B_x}s^2,
\]
and
\[
 \mu_x=\frac{W_x}{2},\qquad \sigma_x=\frac{V_x}{2}.
\]
Let $(\xi_s)_{s\in\B_x}$ be independent Bernoulli variables taking the values
$0,1$ with equal probability, and put $Y_x=\sum_{s\in\B_x}s\xi_s$.

\begin{theorem}[Local limit theorem]\label{thm:lclt}
As $x\to\infty$ through the positive integers,
\begin{equation}\label{eq:lclt}
 \Prob(Y_x=n)=\frac{1}{\sqrt{2\pi}\,\sigma_x}
 \exp\!\left(-\frac{(n-\mu_x)^2}{2\sigma_x^2}\right)
 +o\!\left(\frac1{\sigma_x}\right),
\end{equation}
uniformly for $n\in\Nzero$.  In particular, for all sufficiently large $x$,
every integer $n$ with $|n-\mu_x|\le\sigma_x$ is a subset sum of $\B_x$.
\end{theorem}

\begin{theorem}[Erd\H{o}s Problem~123]\label{thm:main}
For every real $1<\rho<d$, every sufficiently large integer $N$ can be written
as
\[
 N=\sum_{s\in A}s,
 \qquad A\subset \Sg\cap[x,\rho x)
\]
for some positive integer $x$.  Consequently no two distinct elements of $A$
divide one another.
\end{theorem}

The last assertion is immediate.  If $u<v$ belong to $\Sg$ and $u\mid v$,
unique factorization in the three pairwise-coprime bases gives
$v/u=a^r b^s c^t\ge d$; two elements of $[x,\rho x)$ have ratio less than
$\rho<d$.

We prove Theorem~\ref{thm:lclt} for rational $\rho$, then obtain the theorem for
an arbitrary real $\rho$ by choosing a rational $\rho'$ with
$1<\rho'<\rho$.

\section{A triangular grid in the exponent slab}

Put
\[
 \alpha=\log a,\qquad \beta=\log b,\qquad \gamma=\log c,
 \qquad \eta=\log\rho,
\]
and, with $L=\log x$,
\[
 \Eset_L=\{(k,\ell,m)\in\Nzero^3:
 L\le \alpha k+\beta\ell+\gamma m<L+\eta\}.
\]
This set parametrizes $\B_x$.  The ratios of any two of
$\alpha,\beta,\gamma$ are irrational: for instance
$\alpha/\beta=r/s\in\mathbb Q$ would give $a^s=b^r$, contrary to
$\gcd(a,b)=1$.

We use the following elementary consequence of irrational rotation.  If
$u,v>0$, $u/v\notin\mathbb Q$, and $0<\eta<v$, then there is $R$ such that,
whenever
\[
 uX+vY=A+\frac\eta2,\qquad X,Y\ge R,
\]
there are nonnegative integers $r,s$, each within $R$ of the corresponding
real coordinate, for which
\begin{equation}\label{eq:strip-rounding}
 A\le ur+vs<A+\eta.
\end{equation}
Indeed, bounded multiples of $u/v$ form a finite $\eta/(8v)$-net modulo $1$.
Writing $r_0=\lfloor X\rfloor$ and $y=(A-ur_0)/v$, choose a bounded integer
$h$ so that the fractional part of $y-hu/v$ lies between
$1-5\eta/(8v)$ and $1-3\eta/(8v)$; then take
$r=r_0+h$ and $s=\lceil y-hu/v\rceil$.  This gives
\eqref{eq:strip-rounding}, and the displayed linear relation bounds the error
in the second coordinate.

Applying this to the three pairs of logarithms gives a constant $R_0$ with the
following property: every point $(X,Y,Z)$ on
\begin{equation}\label{eq:central-plane}
 \alpha X+\beta Y+\gamma Z=L+\frac\eta2
\end{equation}
whose positive coordinates are at least $2R_0$, and of which at most one is
zero, lies within $R_0$ coordinatewise of a point of $\Eset_L$; a zero
coordinate remains zero.  When all coordinates are positive, round $Y$ to a
nearest integer $\ell_0$, replace
$Z$ by $Z+(\beta/\gamma)(Y-\ell_0)$, and apply
\eqref{eq:strip-rounding} to $\alpha,\gamma$.  If one coordinate is zero,
apply the two-variable statement directly to the other two.

\begin{proposition}[Band geometry]\label{prop:grid}
There are constants $c,C>0$ and integers $D,K\ge1$ such that, for all large
$x$,
\begin{equation}\label{eq:moments}
 cL^2\le |\B_x|\le CL^2,
 \qquad c xL\le V_x\le CxL,
 \qquad c xL^2\le W_x\le CxL^2.
\end{equation}
Moreover, for some $n\asymp L$ there are distinct elements $s_v\in\B_x$
indexed by
\[
 \Delta_n^\circ=\{(i,j,r)\in\Nzero^3:i+j+r=n\}
 \setminus\{(n,0,0),(0,n,0),(0,0,n)\},
\]
such that a zero coordinate of $v$ gives a zero corresponding exponent of
$s_v$, and adjacent vertices satisfy
\begin{equation}\label{eq:edge-relation}
 A_{vw}s_v=B_{vw}s_w,
 \qquad 1\le A_{vw},B_{vw}\le K.
\end{equation}
\end{proposition}

\begin{proof}
Choose $\tau_0$ so large that
$\tau_0/\max(\alpha,\beta,\gamma)>4R_0$.  For large $L$, let
\[
 n=\left\lfloor\frac{L+\eta/2}{\tau_0}\right\rfloor,
 \qquad \tau=\frac{L+\eta/2}{n};
\]
then $n\asymp L$ and $\tau_0\le\tau\le2\tau_0$.  For
$v=(i,j,r)\in\Delta_n^\circ$, round the point
\[
 \left(\frac{\tau i}{\alpha},
       \frac{\tau j}{\beta},
       \frac{\tau r}{\gamma}\right)
\]
to an exponent triple in $\Eset_L$.  The corners were removed, so the rounding
lemma applies.  It preserves the coordinate faces.  Distinct vertices have
one target coordinate separated by more than $4R_0$, hence have distinct
images.  Adjacent targets differ by a bounded amount, so their rounded
exponents differ coordinatewise by at most some fixed $D$.  Clearing the
positive and negative exponent differences gives \eqref{eq:edge-relation} with
$K=(abc)^D$.

Since $|\Delta_n^\circ|\asymp n^2$, this gives $|\B_x|\gg L^2$.  Conversely,
$k,\ell\ll L$, while for fixed $k,\ell$ there is at most one $m$, because two
choices would change the logarithmic weight by at least $\gamma>\eta$.
Thus $|\B_x|\ll L^2$.  Finally $x\le s<\rho x$ throughout the band, and
\eqref{eq:moments} follows by summing $s$ and $s^2$.
\end{proof}

\section{Low-energy frequencies}

For $t\in\T=\R/\Z$, put
\[
 Q_x(t)=\sum_{s\in\B_x}\distz{st}^{2},
 \qquad \distz y=\min_{m\in\Z}|y-m|.
\]
The following is the arithmetic core of the proof.

\begin{proposition}[Rigidity]\label{prop:rigidity}
There are constants $z_0,C,\kappa,\delta>0$ such that, for all large $x$ and
$0\le z\le z_0L^2$,
\begin{equation}\label{eq:low-energy}
 \operatorname{meas}\{t\in\T:Q_x(t)\le z\}
 \le \frac1x
 \exp\!\left(C\left(1+\frac zL\right)\log(L+2)\right).
\end{equation}
The constants may be chosen so that $\delta d\le1/8$ and
\begin{equation}\label{eq:very-low}
 Q_x(t)<\kappa L\quad\Longrightarrow\quad
 \distz t\le\frac\delta x.
\end{equation}
\end{proposition}

\begin{proof}
Use the grid from Proposition~\ref{prop:grid}, and choose
$0<\delta_0<\min(1/(4K),1/32)$.  For each vertex let $m_v$ be a nearest integer to
$s_vt$, and call $v$ bad if $|s_vt-m_v|>\delta_0$.  If there are $H$ bad
vertices, then
\begin{equation}\label{eq:H-bound}
 H\delta_0^2\le Q_x(t).
\end{equation}

For $n/3\le q\le2n/3$, consider the row
\[
 R_q=\{(i,q-i,n-q):0\le i\le q\}.
\]
There are $\gg n$ disjoint such rows, so one contains at most $4H/n$ bad
vertices.  From its central half consider the disjoint paths
\[
 P_j=\{(q-j+r,j,n-q-r):0\le r\le n-q\},
 \qquad q/4\le j\le3q/4.
\]
One of them contains at most $7H/n$ bad vertices.  Their union is connected,
has three endpoints on the three coordinate faces, and contains only
$O(1+H/n)$ bad vertices.  If $8H<n$, the row and path can both be chosen with
no bad vertices.

Orient the edges of this row and path.  For an edge $vw$, define the integer
\[
 e_{vw}=B_{vw}m_w-A_{vw}m_v,
\]
where \eqref{eq:edge-relation} is used.  If both endpoints are good, then
\[
 |e_{vw}|\le B_{vw}|m_w-s_wt|+A_{vw}|m_v-s_vt|<2K\delta_0<1,
\]
so $e_{vw}=0$.  Thus, when $Q_x(t)\le z$, only
$O(1+z/L)$ edge values can be nonzero, by \eqref{eq:H-bound} and $n\asymp L$.
The row, path, positions of the nonzero edge values, and their bounded integer
values can therefore take at most
\begin{equation}\label{eq:code-count}
 \exp\!\left(C\left(1+\frac zL\right)\log(L+2)\right)
\end{equation}
possible values.

Fix these data and compare two compatible frequencies.  If
$\Delta_v$ denotes the difference of their nearest integers, then
\[
 B_{vw}\Delta_w=A_{vw}\Delta_v,
 \qquad B_{vw}s_w=A_{vw}s_v.
\]
Connectivity gives $\Delta_v=r s_v$ for one rational number $r$.  The three
end weights respectively omit $a,b,c$; pairwise coprimality therefore gives
their gcd, and hence the gcd of all weights on the row and path, as $1$.
Consequently $r\in\Z$.  At a fixed root of weight $s_0$, the nearest integer
lies between $0$ and $s_0$, so a fixed choice of the edge data permits at most
three root values modulo $s_0$.  Each root value confines $t$ to an interval of
length $1/s_0\le1/x$.  Multiplying this by \eqref{eq:code-count} proves
\eqref{eq:low-energy}.

If $Q_x(t)<\kappa L$ and $\kappa$ is sufficiently small, then
\eqref{eq:H-bound} gives $8H<n$.  Choose a completely good row and path.  Every
$e_{vw}$ is zero, so $m_v/s_v$ is constant.  Its denominator divides the gcd
of the three face weights, which is $1$; hence the common ratio is an integer.
At any vertex, $|s_vt-m_v|\le\delta_0$ and $s_v\ge x$, so initially
$\distz t\le\delta_0/x$.  To sharpen the constant, apply this conclusion to a
fixed rational subband of width at most $3/2$.  Writing $t=r+u$, we then have
$|u|\le1/(32x)$; nearest-integer rounding vanishes on that subband, and hence
$Q_x(t)\ge u^2\sum s^2\gg x^2L^2u^2$.  Thus $Q_x(t)<\kappa L$ forces
$|u|\le\delta/x$ for any fixed $\delta>0$ once $x$ is large.  Choose
$\delta d\le1/8$ to obtain \eqref{eq:very-low}.
\end{proof}

\section{Fourier analysis and completion of the proof}

We shall use
\begin{equation}\label{eq:cos-decay}
 |\cos(\pi y)|\le \exp(-2\distz y^2),
\end{equation}
which follows from
$\sin(\pi u/2)\ge u$ for $0\le u\le1/2$.

\begin{lemma}[Fourier tails]\label{lem:tails}
\begin{equation}\label{eq:tail}
 \int_{T/V_x\le |t|\le1/2}
 \left|\prod_{s\in\B_x}\cos(\pi st)\right|dt
 =o\!\left(\frac1{V_x}\right),
 \qquad T=L^{1/4}.
\end{equation}
\end{lemma}

\begin{proof}
For large $x$, $T/V_x<2\delta/x$.  On
$T/V_x\le |t|\le2\delta/x$, every $s\in\B_x$ satisfies
$|st|<2\delta d\le1/4$, so \eqref{eq:cos-decay} and
$V_x^2=\sum s^2$ give
\[
 \prod_{s\in\B_x}|\cos(\pi st)|\le \exp(-2V_x^2t^2).
\]
Its integral is at most
$V_x^{-1}\int_T^\infty e^{-2u^2}\,du=o(1/V_x)$.

On $2\delta/x\le |t|\le1/2$, Proposition~\ref{prop:rigidity} gives
$Q_x(t)\ge\kappa L$.  By \eqref{eq:cos-decay}, the integrand is at most
$e^{-2Q_x(t)}$.  Layering by the values of $Q_x$ and using
\eqref{eq:low-energy} yields
\[
 \int e^{-2Q_x(t)}dt
 \le \frac2x\int_{\kappa L}^{z_0L^2}
 e^{-2z+C(1+z/L)\log(L+2)}\,dz+e^{-2z_0L^2}.
\]
For large $L$ this is $O(x^{-1}L^C e^{-3\kappa L/2})+e^{-2z_0L^2}$,
which is $o(1/(xL))=o(1/V_x)$ by \eqref{eq:moments}.
\end{proof}

\begin{proof}[Proof of Theorem~\ref{thm:lclt}]
Fourier inversion gives
\begin{equation}\label{eq:fourier}
 \Prob(Y_x=n)=\int_{-1/2}^{1/2}
 e^{-2\pi i(n-\mu_x)t}
 \prod_{s\in\B_x}\cos(\pi st)\,dt.
\end{equation}
On $|t|\le T/V_x$, all $|\pi st|$ tend uniformly to zero.  The elementary
estimates
\[
 |\cos u-e^{-u^2/2}|\le u^4,
 \qquad
 \left|\prod_j x_j-\prod_j y_j\right|
 \le\sum_j|x_j-y_j|
 \quad (|x_j|,|y_j|\le1)
\]
give
\begin{equation}\label{eq:product-gaussian}
 \left|\prod_{s\in\B_x}\cos(\pi st)
 -\exp\!\left(-\frac{\pi^2V_x^2t^2}{2}\right)\right|
 \le \pi^4t^4\sum_{s\in\B_x}s^4.
\end{equation}
Since $s<px$ and $\sum s^2=V_x^2$,
$\sum s^4\le (px)^2V_x^2$.  Integrating the right side of
\eqref{eq:product-gaussian} over $|t|\le T/V_x$ gives
\[
 O\!\left(\frac{(px)^2T^5}{V_x^3}\right)=o\!\left(\frac1{V_x}\right)
\]
by \eqref{eq:moments}.  The Gaussian tail outside this interval is also
$o(1/V_x)$.  Hence the principal part of \eqref{eq:fourier} equals
\[
 \int_{\R}e^{-2\pi i(n-\mu_x)t}
 \exp\!\left(-\frac{\pi^2V_x^2t^2}{2}\right)dt
 +o\!\left(\frac1{V_x}\right).
\]
The standard Gaussian Fourier transform evaluates this integral as
\[
 \frac{\sqrt{2/\pi}}{V_x}
 \exp\!\left(-\frac{2(n-\mu_x)^2}{V_x^2}\right)
 =\frac1{\sqrt{2\pi}\,\sigma_x}
 \exp\!\left(-\frac{(n-\mu_x)^2}{2\sigma_x^2}\right).
\]
The error is uniform in $n$, since the omitted oscillatory factor always has
modulus $1$.  Lemma~\ref{lem:tails} proves \eqref{eq:lclt}.

If $|n-\mu_x|\le\sigma_x$, the main term in \eqref{eq:lclt} is at least
$e^{-1/2}/(\sqrt{2\pi}\sigma_x)$, whereas the error is
$o(1/\sigma_x)$.  Thus $\Prob(Y_x=n)>0$ for all sufficiently large $x$, so
$n$ is a subset sum of $\B_x$.
\end{proof}

\begin{proof}[Proof of Theorem~\ref{thm:main}]
First suppose $\rho=p/q$ is rational.  By Proposition~\ref{prop:grid}, for all
large $x$ the band has at least $4p^4$ elements, and hence
\begin{equation}\label{eq:V-sweep}
 V_x^2\ge4p^4x^2,
 \qquad V_x\ge2p^2x.
\end{equation}
Also
\begin{equation}\label{eq:W-step}
 W_{x+1}\le W_x+p^2(x+1).
\end{equation}
Indeed, an element entering when $x$ is replaced by $x+1$ lies in an interval
of length $p/q<p$, so there are at most $p$ such integers, each at most
$p(x+1)$; elements leaving the band only decrease the sum.

Fix a sufficiently large initial $X$.  For $N$ large enough that $W_X\le2N$,
choose the largest $x\ge X$ with $W_x\le2N$.  Such an $x$ exists and is finite,
since every sufficiently large band is nonempty and then $W_x\ge x$.  By
maximality and \eqref{eq:W-step},
\[
 0\le2N-W_x<W_{x+1}-W_x\le p^2(x+1)\le2p^2x\le V_x.
\]
Thus $|N-\mu_x|\le\sigma_x$, and Theorem~\ref{thm:lclt} gives a subset of
$\B_x$ summing to $N$.

For a real $1<\rho<d$, choose a rational $\rho'$ with
$1<\rho'<\rho$.  The rational case supplies the representation inside
$[x,\rho'x)\subset[x,\rho x)$.  As observed at the start, every such subset is
a divisibility antichain.
\end{proof}


\end{document}
